\section{Cùng recipe, ViT thua ở đâu}
\label{sec:ch11-cung-recipe}

\index{Vision Transformer!CIFAR-10}%
Chương~\ref{ch:bias-quy-nap} đã có một CIFAR-10 sweep. Tôi không chạy lại ba
mô hình cũ để có một bảng mới đẹp mắt: code và recipe của CNN cùng hai MLP
không đổi, nên output canonical ở tag `ch10`, vẫn có mặt ở tag `ch11`, đã là
hàng tham chiếu. Chương này chỉ chạy 12 mô hình mới: ViT, bốn lượng dữ liệu,
ba seed.

ViT nhỏ này có 66,095 tham số, CNN có 66,570. Nó cắt ảnh thành 16 patch 8x8,
thêm một \tn{token lớp}{class token}, dùng một encoder layer post-LN và một classifier tuyến
tính. Không có convolution, local attention, augmentation hay một thủ thuật
thị giác nào khác. Code, test và output gốc ở tag `ch11`.

\begin{measured}{một ViT cạnh ba hàng của chương 10}
\begin{minted}{text}
n_train  CNN     ViT     MLP matched  MLP wide
1000     0.4365  0.2948  0.3163       0.3469
4000     0.5499  0.3578  0.3733       0.3915
16000    0.6399  0.4150  0.4251       0.4534
50000    0.7036  0.4576  0.4478       0.5025
\end{minted}

Ba seed của ViT có spread lần lượt 0.0274, 0.0094, 0.0210 và 0.0190. CNN có
spread 0.0326, 0.0102, 0.0197 và 0.0019. Chạy bằng
\code{python verify.py --only ch11} tại tag \code{ch11}; gate in
\code{verify: ok}.
\end{measured}

Khoảng cách CNN-ViT là 0.1417 ở 1,000 ảnh và 0.2460 ở 50,000. Cả hai rộng hơn
spread lớn hơn của hàng tương ứng, nên đây không phải một thứ tự do seed lật
được. Nhưng cách đọc quan trọng hơn là hàng ViT cạnh MLP: ở 50,000 ảnh, ViT
vẫn thua CNN mà vượt MLP khớp tham số một ít. Bỏ cấu trúc cục bộ không đồng
nghĩa với trở thành MLP; attention vẫn cho patch nói chuyện toàn cục từ layer
đầu.

\subsection*{Nó chưa khớp cả phần huấn luyện được chấm}

Chỉ nhìn test accuracy thì dễ kể sai câu chuyện: ViT tổng quát hóa kém. Cột
train accuracy chặn câu đó.

\begin{measured}{độ chính xác trên tối đa 10,000 ảnh huấn luyện đầu}
\begin{minted}{text}
n_train  ViT train acc
1000     0.5297
4000     0.4306
16000    0.4418
50000    0.4644
\end{minted}

Đo bằng \code{experiments/ch11\_vit.py} tại tag \code{ch11}.
\end{measured}

Hàm chấm của companion repo cố ý chỉ đọc tối đa 10,000 ảnh huấn luyện, để không
đốt thêm thời gian cho một cột chẩn đoán. Vì vậy ở hai hàng 16,000 và 50,000,
đây không phải accuracy trên toàn bộ tập train. Dù với giới hạn ấy, ViT chỉ
khớp 0.5297 ở hàng dễ nhất, rồi còn thấp hơn khi dữ liệu nhiều lên vì số epoch
giảm để giữ số bước gần cùng bậc với chương 10. Nên bảng này không được phép
nói \enquote{ViT overfit trên CIFAR-10}. Nó chưa khớp cả phần dữ liệu huấn luyện
được chấm trong ngân sách recipe này. Đó là một kết quả tệ hơn cho ViT, nhưng
sạch hơn: không phải một câu chuyện sau khi mô hình đã học xong.

Tôi không nâng số epoch cho ViT để nó đẹp hơn. Làm thế sẽ đổi hai biến cùng
lúc, kiến trúc và ngân sách huấn luyện, rồi không còn biết số nào trả tiền cho
số nào.
